% Описание_назначения_документов.tex — DOC-01 v1.0
\documentclass[12pt,a4paper]{article}

% =============================================================================
% КОНФИГУРАЦИЯ И ПРЕАМБУЛА
% =============================================================================

% --- конфигурация метаданных ---
\ProvidesFile{config.tex}[2026/10/24 mephi-12 config]

% ---------- Прямые значения (редактируйте по необходимости) ----------
\newcommand{\Company}{НИЯУ МИФИ}
\newcommand{\DocumentTitleMain}{Описание назначения документов}
\newcommand{\DocumentTitleSub}{комплект документации криптотренажёра mephi-12}
\newcommand{\Version}{1.0}
\newcommand{\DocDate}{24.10.2026}
\newcommand{\AuthorName}{Самбьюи А.Ж.}
\newcommand{\AgreedBy}{}
\newcommand{\ApprovedBy}{}
\newcommand{\City}{Москва}

% --- преамбула и стили ---
\usepackage[utf8]{inputenc}
\usepackage[T2A]{fontenc}
\usepackage[russian]{babel}

% Times New Roman-like
\usepackage{newtxtext,newtxmath} % похож на Times; обеспечивает текст+математику
\usepackage{setspace}
\onehalfspacing % межстрочный интервал 1.5

% Страницы
\usepackage{geometry}
\geometry{left=25mm,right=15mm,top=20mm,bottom=20mm}

% отступ первой строки
\usepackage{indentfirst}
\setlength{\parindent}{1.25cm}

% таблицы, картинки, списки, подписи
\usepackage{graphicx}
\usepackage{caption}
\usepackage{subcaption}
\captionsetup[table]{position=top}    % подпись таблицы над ней
\captionsetup[figure]{position=bottom}% подпись рисунка под ней
\usepackage{booktabs}
\usepackage{tabularx}
\usepackage{longtable}
\usepackage{array}
\usepackage{enumitem}
\setlist[itemize]{left=0pt,label=\textbullet}
\setlist[enumerate]{left=0pt,label=\arabic*.}

% код и конфигурации
\usepackage{listings}
\lstset{basicstyle=\ttfamily\small,breaklines=true}

% гиперссылки
\usepackage{hyperref}
\hypersetup{colorlinks=true,linkcolor=black,urlcolor=blue}

\pagestyle{plain}  % только номер страницы внизу по центру

% нумерация разделов — стандартная десятичная (article по умолчанию)
% дополнительные удобства
\usepackage{tocloft} % управление оглавлением
\setlength{\cftbeforesecskip}{6pt}

% небольшие утилиты
\newcommand{\reqfield}[1]{\textbf{#1}:} % для шаблонов в приложениях

\begin{document}

% =============================================================================
% ТИТУЛЬНЫЙ ЛИСТ
% =============================================================================
\begin{titlepage}
  \thispagestyle{empty} % убирает номер/колонтитул на этой странице
  \centering
  \vspace*{5.8cm}
  {\Huge \textbf{\DocumentTitleMain}\par}
  {\Large \vspace{4pt} \DocumentTitleSub \par} % вторая часть — обычным шрифтом, чуть ниже
  \vspace{0.8cm}

  {\large Версия: \Version \par}
  {\large Дата: \DocDate \par}
  \vspace{1.2cm}

  % ----- Таблица с ролями -----
  \centering
  \begin{tabularx}{0.95\textwidth}{@{\hspace{6pt}}p{5.5cm} p{6cm} p{4cm}@{\hspace{6pt}}}
    \toprule
    \textbf{Роль} & \textbf{ФИО} & \textbf{Подпись / Дата} \\
    \midrule
    Разработал: & \AuthorName & \\
    \addlinespace[6pt]
    Согласовано: & \AgreedBy & \\
    \addlinespace[6pt]
    Утверждено: & \ApprovedBy & \\
    \bottomrule
  \end{tabularx}

  \vfill
  {\large г. \City, \the\year \par}
\end{titlepage}

% =============================================================================
% ЖУРНАЛ ИЗМЕНЕНИЙ
% =============================================================================
\section*{Лист регистрации изменений \texttt{[DPURP.CHANGES]}}
\addcontentsline{toc}{section}{Лист регистрации изменений}

\begin{table}[h!]
\centering
\small
\begin{tabularx}{\textwidth}{
  >{\centering\arraybackslash}p{1.1cm}
  >{\centering\arraybackslash}p{2.1cm}
  p{2.3cm}
  p{2.6cm}
  p{1.6cm}
  X
  p{1.8cm}
}
\toprule
\textbf{Версия} & \textbf{Дата} & \textbf{Автор} & \textbf{Изменённый элемент} & \textbf{Х-р изменения} & \textbf{Описание изменения} & \textbf{Утвердил} \\
\midrule
1.0 & 24.10.2026 & Самбьюи А.Ж. & Весь документ & Добавлено & Первичный выпуск документа. & \\
\bottomrule
\end{tabularx}
\end{table}

\clearpage
\tableofcontents
\clearpage

% =============================================================================
% ОСНОВНОЕ СОДЕРЖАНИЕ
% =============================================================================

\section{Введение \texttt{[DPURP.INTRO]}}

\subsection{Назначение документа \texttt{[DPURP.INTRO.PURPOSE]}}
Настоящий документ является путеводителем по комплекту проектной документации
криптотренажёра mephi-12: он определяет, какие документы входят в комплект, какую
роль каждый из них играет в жизненном цикле проекта, кто является его основным
потребителем и на каких входных материалах он основывается. Документ помогает
читателю с любой ролью (заказчик, разработчик, рецензент, верификатор) быстро
определить, какие документы ему необходимы и в каком порядке их изучать.
Перечень документов, статусы версий и календарь выпуска приведены в README
комплекта (DOC-00); настоящий документ дополняет его сведениями о назначении
и взаимосвязях.

\subsection{Наименование системы \texttt{[DPURP.INTRO.NAME]}}
Объектом документирования является программная система с полным наименованием
«Платформа для изучения и практики курса по криптографии в МИФИ» и кратким
наименованием \textbf{криптотренажёр mephi-12}. Система реализована как backend-монолит
на Kotlin 1.9.25 + Spring Boot 3.4.0 с СУБД PostgreSQL 17, аутентификацией по JWT
и развёртыванием в Docker; включает 7 REST-контроллеров (14 действующих эндпоинтов,
базовый путь \texttt{/api}) и 8 JPA-сущностей. Поддерживаемые типы задач:
ранцевые криптосистемы (супервозрастающая и на степенях), атака бэкдором на
ранцевую криптосистему и схема Эль-Гамаля.

\subsection{Стандарты-ориентиры \texttt{[DPURP.INTRO.STD]}}
Комплект составлен по образцу авиационных стандартов \textbf{DO-178C} и его
российского аналога \textbf{КТ-178C} (ссылочные издания — в папке
\texttt{00\_Справочные\_стандарты\_и\_методички}). Из стандартов заимствованы:
каскад «техническое задание $\rightarrow$ системные требования $\rightarrow$
программные требования $\rightarrow$ трассировка», формальные инспекции по методу
Фагана, планы верификации и управления конфигурацией ПО, а также сквозная
система идентификаторов требований (\texttt{[SR.*]}, \texttt{[SWR.*]}, TC-XXX).
Комплект является учебной адаптацией стандартов в рамках курсовой работы НИЯУ МИФИ:
тестовый код и прогоны тестов не выполняются (решение от 20.10), верификация
представлена планами с описательными кейсами TC-XXX.

\section{Категории потребителей \texttt{[DPURP.AUDIENCE]}}
\begin{itemize}
  \item \textbf{Заказчик / руководитель курса.} Оценивает полноту и качество комплекта
        в целом. Ему предназначены DOC-00 (README комплекта — карта репозитория),
        DOC-01 (настоящий документ), DOC-02 (что заказывается), DOC-09 и DOC-14
        (доказательство покрытия требований), а также протоколы формальных инспекций
        DOC-13.1--13.6 как свидетельство контроля качества.
  \item \textbf{Разработчики — авторы документов.} Формируют содержание комплекта.
        Работают по ПУКПО (DOC-03: элементы конфигурации, ветки, теги), пишут
        требования в DOC-05 и DOC-06.1--06.6 по своим доменам, оформляют архитектуру
        DOC-07 и справочник API DOC-08; опираются на методику идентификаторов и
        глоссарий (материалы \texttt{notes/methodology}).
  \item \textbf{Рецензенты / инспекторы.} Проводят формальные инспекции по процедуре
        DOC-04 (роли: автор, ведущий, секретарь, рецензенты; правило перекрёстности —
        рецензент проверяет чужой домен). Предмет инспекции — проверяемый документ
        (DOC-02, DOC-03, DOC-05, DOC-06, DOC-07/08, DOC-09/10/11); результат
        оформляется протоколами DOC-13.1--13.6 и учитывается в реестре дефектов
        \texttt{defects.md}.
  \item \textbf{Верификаторы.} Подтверждают, что требования прослеживаются до
        верификационных кейсов. Используют DOC-10 (план верификации, кейсы TC-XXX),
        DOC-11 (план интеграционного тестирования), DOC-09 (матрица трассировки
        SR $\rightarrow$ SWR $\rightarrow$ TC) и исходные требования DOC-05, DOC-06.
\end{itemize}

\section{Назначение и статус каждого документа \texttt{[DPURP.DOCS]}}

\subsection{Назначение и зависимости \texttt{[DPURP.DOCS.TABLE]}}
В таблице~\ref{tab:purpose-docs} приведены все документы комплекта DOC-00--DOC-14
(кроме исключённого DOC-12): назначение, основной потребитель, входные зависимости
и приоритет.

{\small
\setlength{\tabcolsep}{3.5pt}
\begin{longtable}{>{\raggedright\arraybackslash}p{1.3cm} >{\raggedright\arraybackslash}p{2.7cm} >{\raggedright\arraybackslash}p{4.7cm} >{\raggedright\arraybackslash}p{2.4cm} >{\raggedright\arraybackslash}p{2.6cm} >{\centering\arraybackslash}p{0.7cm}}
\caption{Назначение, потребители и зависимости документов комплекта}
\label{tab:purpose-docs}\\
\toprule
\textbf{ID} & \textbf{Документ} & \textbf{Назначение} & \textbf{Потребитель} & \textbf{Зависимости} & \textbf{Приор.} \\
\midrule
\endfirsthead
\toprule
\textbf{ID} & \textbf{Документ} & \textbf{Назначение} & \textbf{Потребитель} & \textbf{Зависимости} & \textbf{Приор.} \\
\midrule
\endhead
\bottomrule
\endlastfoot
DOC-00 & README комплекта & Входная точка: оглавление, структура папок и перечень всех документов с версиями. & Все категории & Структура репозитория & К \\
DOC-01 & Описание назначения документов & Путеводитель по комплекту: назначение, потребитель и зависимости каждого документа (настоящий документ). & Заказчик, новые участники & DOC-00 & В \\
DOC-02 & Техническое задание \texttt{[TS.*]} & Контракт на разработку: цели, границы и требования заказчика к платформе. & Заказчик, разработчики & Карточки доменов (все) & К \\
DOC-03 & ПУКПО \texttt{[SCMP.*]} & Правила управления конфигурацией: элементы конфигурации, версии, ветки и теги, роли. & Разработчики, секретарь ФИ & Репозиторий, роли & К \\
DOC-04 & Процедура формальной инспекции & Регламент инспекций по Фагану: этапы, роли, чек-листы, категории замечаний. & Инспекторы (все роли ФИ) & По образцу & К \\
DOC-05 & Системные требования \texttt{[SR.*]} & Что система должна делать с точки зрения заказчика; источник для SWR и трассировки. & Разработчики, верификаторы & DOC-02, методика ID & К \\
DOC-06 & Программные требования, сводные & Распределение \texttt{[SWR.*]} по подсистемам, архитектура ПО, модель данных (8 сущностей). & Разработчики, верификаторы & DOC-05 & К \\
DOC-06.1 & SWR: аутентификация и пользователи & Программные требования подсистемы регистрации, JWT и управления пользователями. & Разработчики, верификаторы & DOC-05 & К \\
DOC-06.2 & SWR: сессии задач и планировщик & Программные требования подсистемы сессий задач и автозавершения по таймауту. & Разработчики, верификаторы & DOC-05 & К \\
DOC-06.3 & SWR: ранцевые КС & Программные требования генерации и решения ранцевых криптосистем двух видов. & Разработчики, верификаторы & DOC-05 & К \\
DOC-06.4 & SWR: бэкдор в ранцевой КС & Программные требования подсистемы демонстрации атаки бэкдором. & Разработчики, верификаторы & DOC-05 & В \\
DOC-06.5 & SWR: Эль-Гамаль и матбиблиотека & Программные требования схемы Эль-Гамаля и утилит математической библиотеки. & Разработчики, верификаторы & DOC-05 & К \\
DOC-06.6 & SWR: конфигурации и модерация & Программные требования подсистемы конфигураций задач и модерации. & Разработчики, верификаторы & DOC-05 & В \\
DOC-07 & Архитектура & Компоненты, слои и схема развёртывания ПО; связывает SWR с реализацией. & Разработчики & Обзор кода, DOC-06 & В \\
DOC-08 & API — справочник & Справочник всех REST-эндпоинтов (14 действующих): методы, пути, роли, контракты. & Разработчики, верификаторы & Контроллеры & В \\
DOC-09 & Матрица трассировки & Доказательство покрытия: каждое SR связано с SWR и верификационным кейсом TC. & Заказчик, верификаторы & DOC-05, DOC-06, DOC-10 & К \\
DOC-10 & План верификации ПО & Как подтверждается выполнение каждого требования (кейсы TC-XXX). & Верификаторы & ID требований СТ/SWR & К \\
DOC-11 & План интеграционного тестирования & Сценарии проверки взаимодействия подсистем через REST API. & Верификаторы & DOC-06, DOC-10 & В \\
DOC-13.1 & Протокол ФИ-ТЗ & Фиксирует результаты формальной инспекции технического задания. & Заказчик, автор DOC-02 & DOC-04, DOC-02 v1.0 & К \\
DOC-13.2 & Протокол ФИ-ПУКПО & Фиксирует результаты формальной инспекции ПУКПО. & Заказчик, автор DOC-03 & DOC-04, DOC-03 v1.0 & К \\
DOC-13.3 & Протокол ФИ-СТ & Фиксирует результаты формальной инспекции системных требований. & Заказчик, автор DOC-05 & DOC-04, DOC-05 v1.0 & К \\
DOC-13.4 & Протокол ФИ-SWR & Фиксирует результаты формальной инспекции программных требований. & Заказчик, авторы DOC-06 & DOC-04, DOC-06 v1.0 & К \\
DOC-13.5 & Протокол ФИ-АРХ+API & Фиксирует результаты формальной инспекции архитектуры и API. & Заказчик, авторы DOC-07/08 & DOC-04, DOC-07/08 v1.0 & К \\
DOC-13.6 & Протокол ФИ-TRACE+VER & Фиксирует результаты инспекции трассировки и планов верификации. & Заказчик, авторы DOC-09/10/11 & DOC-04, DOC-09/10/11 & К \\
DOC-14 & Анализ покрытия (RCOV) & Количественная оценка покрытия требований для приёмки комплекта. & Заказчик, руководитель курса & DOC-09 v1.0 & О \\
\end{longtable}
}

\noindent Примечания.
\begin{enumerate}
  \item Приоритет: \textbf{К} — критический путь (задержка документа сдвигает весь
        комплект), \textbf{В} — высокий, \textbf{О} — опционально (выполняется при
        наличии ресурса; комплект полон и без него).
  \item DOC-12 (протоколы прогонов тестов TPR) исключён из комплекта решением
        от 20.10: тестовый код и прогоны в рамках работы не выполняются.
  \item Трассировка (DOC-09) ведётся в CSV, в комплект включается её PDF-версия.
\end{enumerate}

\subsection{Статус на 24.10.2026 \texttt{[DPURP.DOCS.STATUS]}}
По завершении Спринта~1 выпущены: DOC-00 v1.0, DOC-01 v1.0 (настоящий документ),
DOC-04 v1.0, DOC-02 v0.5, DOC-03 v0.5 и скелет DOC-05 v0.1; подготовлены рабочие
материалы (карточки доменов, каталог эндпоинтов, методика идентификаторов).
Остальные документы выпускаются по графику Спринтов~2--3 и финализации;
актуальные версии и даты приведены в DOC-00.

\section{Жизненный цикл комплекта \texttt{[DPURP.LIFE]}}

\subsection{Этапность по спринтам \texttt{[DPURP.LIFE.SPRINTS]}}
Комплект создаётся за три спринта с последующей финализацией (таблица~\ref{tab:life-sprints}).

\begin{table}[h!]
    \centering
    \caption{Этапы жизненного цикла комплекта}
    \label{tab:life-sprints}
    \begin{tabularx}{\textwidth}{p{3.2cm} p{2.6cm} X}
        \toprule
        \textbf{Этап} & \textbf{Период} & \textbf{Основные результаты} \\
        \midrule
        Спринт 1 — фундамент & 20.10--24.10.2026 & Структура репозитория, шаблон LaTeX, методика ID, карточки доменов, каталог эндпоинтов; DOC-00 v1.0, DOC-01 v1.0, DOC-04 v1.0, DOC-02 v0.5, DOC-03 v0.5, скелет DOC-05 v0.1. \\
        Спринт 2 — СТ и SWR, первые ФИ & 25.10--07.11.2026 & DOC-02 v1.0--v1.1, DOC-03 v1.0, DOC-05 v1.0, DOC-06 v0.5, DOC-06.1--06.6 v0.8, DOC-07/08 v0.5, DOC-11 v0.3, черновик CSV трассировки; ФИ-ТЗ, ФИ-ПУКПО, ФИ-СТ. \\
        Спринт 3 — v1.0 всех документов, трассировка & 08.11--21.11.2026 & DOC-05 v1.1, DOC-06 v1.0 с моделью данных, DOC-06.1--06.6 v1.0, DOC-07/08 v1.0, DOC-09 v1.0, DOC-10 v1.0, DOC-11 v1.0; ФИ-SWR, ФИ-АРХ+API, ФИ-TRACE+VER. \\
        Финализация — правки, сборка, приёмка & 22.11--01.12.2026 & DOC-09 v1.1, DOC-14 v1.0, закрытие замечаний ФИ, сквозной контроль трассировки, титулы и листы изменений, сборка и приёмка комплекта. \\
        \bottomrule
    \end{tabularx}
\end{table}

\subsection{Взаимосвязь документов \texttt{[DPURP.LIFE.FLOW]}}
Документы образуют сквозную цепочку вывода требований и доказательств:

\begin{center}
ТЗ (DOC-02) $\rightarrow$ СТ (DOC-05) $\rightarrow$ SWR (DOC-06, DOC-06.1--06.6) $\rightarrow$ \\
$\rightarrow$ \{Архитектура (DOC-07), API (DOC-08)\} $\rightarrow$ \{Планы верификации (DOC-10, DOC-11)\} $\rightarrow$ \\
$\rightarrow$ Трассировка (DOC-09) $\rightarrow$ ФИ-протоколы (DOC-13.1--13.6).
\end{center}

Техническое задание задаёт цели и границы системы; системные требования детализируют
его в проверяемые формулировки \texttt{[SR.*]}; программные требования распределяют
их по подсистемам \texttt{[SWR.*]} и структурам данных; архитектура и API описывают
реализацию требований; планы верификации определяют кейсы TC-XXX; матрица
трассировки замыкает цепочку, связывая каждое SR с SWR и TC; формальные инспекции
(по процедуре DOC-04) подтверждают качество каждой ступени. Организационный контур
образуют ПУКПО (DOC-03 — версии, ветки, теги) и реестр дефектов; навигационный —
README (DOC-00) и настоящий документ.

\end{document}
